\chapter*{Anhang}
\addcontentsline{toc}{chapter}{Anhang}

\section*{Hinweis zu den Anlagen}
Die folgenden Anlagen sind gem\"a\ss{} Vorgabe beizuf\"ugen. Regeln:
\begin{itemize}
  \item Gesamtumfang aller Anlagen: max.\ 20 Seiten.
  \item Auf jede hier gelistete Anlage muss im Hauptteil eindeutig verwiesen
    werden (per \verb|\ref{}|, siehe Beispiele in den Kapiteln).
  \item Anlagen, die nicht selbst erstellt wurden (offizielle Formulare
    o.\,\"a.), sind als solche zu kennzeichnen -- siehe Anlage~1 und Anlage~18.
  \item Vertrauliche/datenschutzrelevante Inhalte in Anlagen sind ebenso zu
    kennzeichnen wie im Hauptteil.
  \item Nicht ben\"otigte Anlagen (z.\,B. wenn im eigenen Projekt kein
    Gantt-Diagramm separat von der Zeitplanung existiert) l\"oschen und die
    zugeh\"orige \verb|\input|-Zeile unten entfernen -- nicht leer stehen
    lassen.
  \item Vor dem Hochladen pr\"ufen, ob eingebundene PDFs (insbesondere
    gescannte Unterschriften) komprimiert sind (Gesamtgr\"o\ss{}e beachten).
\end{itemize}

\section*{Anlage 1: Projektantrag}
\label{anh:projektantrag}
Der urspr\"unglich bei der IHK eingereichte und genehmigte Projektantrag.
Dient als Nachweis, dass die Durchf\"uhrung erst nach Genehmigung begonnen
wurde, und als Referenz f\"ur den Soll-Ist-Vergleich in
Abschnitt~\ref{sec:soll_ist_vergleich}.

[PLATZHALTER: Projektantrag als PDF einbinden, z.\,B.:]
\begin{verbatim}
% \includepdf[pages=-]{anlagen/projektantrag.pdf}
\end{verbatim}

\section*{Anlage 2: Lastenheft}
\label{anh:lastenheft}
Referenziert in Abschnitt~\ref{sec:lastenheft}.

[PLATZHALTER: Vollst\"andiges Lastenheft (Anforderungen aus Nutzersicht)
einf\"ugen -- entweder als eigenes Kapitel hier oder als eingebundenes PDF.]

\section*{Anlage 3: Pflichtenheft}
\label{anh:pflichtenheft}
Referenziert in Abschnitt~\ref{sec:pflichtenheft}.

[PLATZHALTER: Vollst\"andiges Pflichtenheft (technische Umsetzungsvorgaben)
einf\"ugen.]

\section*{Anlage 4: Gantt-Diagramm}
\label{anh:gantt}
Referenziert in Abschnitt~\ref{sec:zeitplanung}.

[PLATZHALTER: Grafisches Gantt-Diagramm der Projektphasen einf\"ugen, z.\,B.:]
\begin{verbatim}
% \includegraphics[width=\linewidth]{bilder/gantt.png}
\end{verbatim}

\section*{Anlage 5: Zeitplanung}
\label{anh:zeitplanung}
Referenziert in Abschnitt~\ref{sec:zeitplanung}.

[PLATZHALTER: Detaillierte Zeitplanung (Phasen, geplante/tats\"achliche
Stunden) als Tabelle einf\"ugen.]

\begin{table}[h!]
  \centering
  \caption{Zeitplanung}
  \begin{tabular}{@{}p{0.6\linewidth}rr@{}}
    \toprule
    \textbf{Projektphase} & \textbf{Geplant (h)} & \textbf{Ist (h)} \\
    \midrule
    {[}PLATZHALTER{]} & {[}PLATZHALTER{]} & {[}PLATZHALTER{]} \\
    \bottomrule
  \end{tabular}
\end{table}

\section*{Anlage 6: Kostenplanung}
\label{anh:kostenplanung}
Referenziert in Abschnitt~\ref{ch:projektplanung}.

[PLATZHALTER: Kostenaufstellung (Personal-, Sach-, Lizenzkosten) und ggf.
Amortisationsrechnung/Wirtschaftlichkeitsbetrachtung einf\"ugen.]

\section*{Anlage 7: Risikoanalyse}
\label{anh:risikoanalyse}
Referenziert in Abschnitt~\ref{ch:projektplanung}.

[PLATZHALTER: Vollst\"andige Risikoanalyse (Risiko, Eintrittswahrscheinlichkeit,
Auswirkung, Gegenma\ss{}nahme) als Tabelle einf\"ugen.]

\section*{Anlage 8: Stakeholderanalyse}
\label{anh:stakeholderanalyse}
Referenziert in Abschnitt~\ref{ch:projektplanung}.

[PLATZHALTER: Stakeholdermatrix (Einfluss/Interesse) einf\"ugen.]

\section*{Anlage 9: Vorgehensmodell}
\label{anh:vorgehensmodell}
Referenziert in Abschnitt~\ref{ch:vorgehensmodell}.

[PLATZHALTER: Grafische Darstellung des gew\"ahlten Vorgehensmodells
(z.\,B. Phasenmodell, Sprint-\"Ubersicht) einf\"ugen.]

\section*{Anlage 10: Testprotokoll}
\label{anh:testprotokoll}
Referenziert in Kapitel~\ref{ch:testphase_und_fehlerbehebung}.

[PLATZHALTER: Vollst\"andiges Testprotokoll (Testfall, erwartetes Ergebnis,
tats\"achliches Ergebnis, Status) als Tabelle einf\"ugen.]

\section*{Anlage 11: Diagramme (Use-Case, ER-/Tabellenmodell, Aktivit\"atsdiagramme)}
\label{anh:diagramme}
Referenziert in Kapitel~\ref{ch:umsetzung}.

[PLATZHALTER: Use-Case-Diagramm, ER-Modell/Tabellenmodell,
Komponenten-/Klassendiagramm sowie Aktivit\"atsdiagramme der wichtigsten
Abl\"aufe einf\"ugen (je Diagrammtyp eigene \verb|\subsection*|, falls mehr
als eines).]

\section*{Anlage 12: Mockups}
\label{anh:mockups}
Referenziert in Abschnitt~\ref{ch:umsetzung}.

[PLATZHALTER: UI-Entw\"urfe/Mockups vor der Umsetzung einf\"ugen.]

\section*{Anlage 13: Screenshots der Anwendung}
\label{anh:screenshots}
Referenziert in Abschnitt~\ref{ch:umsetzung}.

[PLATZHALTER: Screenshots der fertigen Anwendung einf\"ugen. Bei
Kundendaten/internen Daten in Screenshots: unkenntlich machen oder als
vertraulich kennzeichnen.]

\section*{Anlage 14: Entwicklerdokumentation}
\label{anh:entwicklerdokumentation}
Referenziert in Abschnitt~\ref{sec:soll_ist_vergleich}.

[PLATZHALTER: Technische Dokumentation f\"ur Nachfolger/Wartung
(Installation, Architektur, wichtige Klassen/Module) einf\"ugen.]

\section*{Anlage 15: Kundendokumentation}
\label{anh:kundendokumentation}
Referenziert in Abschnitt~\ref{sec:soll_ist_vergleich}.

[PLATZHALTER: Anleitung f\"ur Endnutzer (Bedienung der Anwendung) einf\"ugen.]

\section*{Anlage 16: Quellcode-Ausschnitte}
\label{anh:quellcode}
Referenziert in Kapitel~\ref{ch:umsetzung}.

[PLATZHALTER: Ausgew\"ahlte, besonders erkl\"arungsbed\"urftige
Quelltext-Ausschnitte einf\"ugen -- nicht der vollst\"andige Quelltext
(Seitenlimit). Z.\,B.:]
\begin{verbatim}
% \lstinputlisting[language=Python, firstline=1, lastline=40]{code/beispiel.py}
\end{verbatim}

\section*{Anlage 17: Protokoll zur Projektarbeit}
\label{anh:protokoll_projektarbeit}
Nicht selbst erstellt -- offizielles, von der IHK bereitgestelltes Formular
(zum Herunterladen auf der IHK-Downloadseite verf\"ugbar). Als
unterschriebenes PDF einbinden:
\begin{verbatim}
% \includepdf[pages=-]{anlagen/protokoll_projektarbeit.pdf}
\end{verbatim}
[PLATZHALTER: Datei ablegen unter \texttt{anlagen/protokoll\_projektarbeit.pdf}
und obige Zeile aktivieren.]

\section*{Anlage 18: Pers\"onliche Erkl\"arung}
\label{anh:persoenliche_erklaerung}
Nicht selbst erstellt -- offizielles, von der IHK bereitgestelltes Formular
(zum Herunterladen auf der IHK-Downloadseite verf\"ugbar). Als
unterschriebenes PDF einbinden:
\begin{verbatim}
% \includepdf[pages=-]{anlagen/persoenliche_erklaerung.pdf}
\end{verbatim}
[PLATZHALTER: Datei ablegen unter \texttt{anlagen/persoenliche\_erklaerung.pdf}
und obige Zeile aktivieren.]

